\section{Related Work}
\label{submission:sec:related}

\textbf{Model Merging in Parameter Space.} Model merging has its roots in early parameter averaging techniques like Federated Averaging (FedAvg)~\cite{mcmahan2017communication} and Model Soups~\cite{wortsman2022model}. FedAvg averages models trained on decentralized partitions, while Model Soups averages multiple fine-tuned models from different hyperparameter configurations to improve generalizability. 
Building upon these concepts, Ilharco et al.~\cite{ilharco2022editing} proposed \textbf{Task Arithmetic}, which defines "task vectors" as the difference in parameters between a fine-tuned model and its pretrained initialization. Task vectors can be added or subtracted to steer the base model's behavior. To improve merging accuracy, several works integrate statistics from the training data: RegMean~\cite{jin2023raw} solves a closed-form least-squares problem using feature covariance matrices, and Fisher Merging~\cite{matena2022merging} computes the empirical Fisher information matrix to weight each parameter. While successful, these approaches require storing and computing auxiliary data statistics, which can be computationally heavy and logistically restrictive when training datasets are unavailable.

\textbf{Resolving Parameter Interference.} A major challenge in parameter merging is interference, which arises when different models have conflicting parameter updates. Yadav et al.~\cite{yadav2023resolving} introduced \textbf{Ties-Merging}, which prunes low-magnitude parameters (Trimming), resolves sign conflicts via voting (Elect Sign), and averages only the non-conflicting updates (Disjoint Merge). Similarly, \textbf{DARE}~\cite{yu24dare} uses a randomized pruning mechanism to sparsify task vectors before merging, compensating for the pruned weights via magnitude scaling. While Ties-Merging and DARE are powerful at mitigating sign and parameter conflicts, they rely on heuristic hyperparameters (such as pruning ratios or drop rates) that must be manually tuned. Furthermore, they do not resolve systematic representation scale mismatches across different tasks, which can lead to a single task dominating the entire merged representation.

\textbf{Optimization-Based and Complex Merging.} 
To avoid manual hyperparameter tuning, recent works have turned to optimization-based or algebraic-heavy procedures during or after merging. \textbf{AdaMerging}~\cite{yang2024adamerging} formulates merging coefficients as learnable parameters, optimization-tuned via active entropy minimization on a small unlabeled test batch. While adaptive, this introduces test-time optimization loops, requires access to unlabeled data, and is susceptible to suboptimal local minima. 

Even more complex pipelines have emerged to balance representations or exploit geometric structure. \textbf{SyMerge}~\cite{symerge2025synergistic} optimizes a task-specific layer at test-time under the guidance of expert teacher models. \textbf{OrthoMerge}~\cite{orthomerge25manf} performs magnitude-corrected merging on the Riemannian manifold of the orthogonal group using singular value decomposition (SVD), Cayley transforms, and Procrustes alignment. \textbf{SAIM}~\cite{saim25flat} utilizes Sharpness-Aware Block Coordinate Descent (SA-BCD) during training, followed by Adaptive Isotropic Merging which scales cumulative updates based on SVD of the parameter matrices. 

Although these methods address representation imbalance, they introduce tremendous computational overhead. SVD scales cubically $O(d^3)$ with weight matrix dimensions, which is prohibitively expensive for large layers, while active test-time optimization incurs significant latency. In contrast, our proposed \textbf{SD-Scale} operates strictly in linear time $O(N)$ (where $N$ is the parameter count) using simple element-wise operations. It is completely training-free, non-parametric, and requires no auxiliary SVDs or unlabeled test data, preserving the minimalist elegance and speed of parameter-space merging.
